The behavior of a language model is shaped not only by its weights but also by
its \emph{harness}: prompts, memory, skills, tools, middleware, and control
flow.  Recent systems automatically rewrite these surfaces, yet an adaptively
selected score increase need not establish that the edit was active, followed,
or causally responsible, and repeated promotion tests can overfit their own
validation process.  We formulate self-evolving harnesses as a verification
problem and introduce \systemname, a training-free, mechanism-gated protocol.
Candidates are typed atomic mutations; a promotion authority outside the
mutable harness evaluates matched child, parent, revert, and neutral-placebo
interventions, then checks activation, adherence, intended behavior, utility,
protected behavior, and cost under a campaign-level error budget.  We also
specify a prospective HarnessFaultBench design with known causal surfaces and a
budget-matched joint-policy comparison between mechanism-visible,
mechanism-gated and score-only, performance-gated search using the same exact
model configuration.  A controlled-fault feedback-by-promotion factorial is
designed to separate conditional policy-component effects.  Confirmatory claims
remain
\evidenceblocked{the protocol is not frozen and the independent searches and
one-time sealed evaluation have not been completed}.  This explicit boundary
prevents development pilots from being mistaken for evidence that mechanism
gating improves self-evolution.
